\chapter{ตัวอย่างการทำงานของระบบ Text-to-SQL}

ภาคผนวกนี้แสดงตัวอย่างคำถามภาษาไทยของผู้ใช้ คู่กับคำสั่ง SQL ที่โมเดลภาษาสร้างขึ้นและนำไปประมวลผลบนฐานข้อมูล เพื่อให้เห็นภาพการทำงานจริงของแกน Text-to-SQL

\vskip0.5cm
\noindent\textbf{ตัวอย่างที่ 1} คำถาม: ``หลักสูตร IT ต้องเรียนกี่หน่วยกิตถึงจะจบ''

\begin{verbatim}
SELECT program_code, total_credits
FROM programs
WHERE program_code = 'IT';
\end{verbatim}

\noindent\textbf{ตัวอย่างที่ 2} คำถาม: ``วิชาบังคับก่อนของวิชารหัส 06016418 มีอะไรบ้าง''

\begin{verbatim}
SELECT cp.prerequisite_code
FROM course_prerequisites cp
JOIN courses c ON c.course_id = cp.course_id
WHERE c.code = '06016418';
\end{verbatim}

\noindent\textbf{ตัวอย่างที่ 3} คำถาม: ``แผนการเรียนปกติของ IT ปี 1 เทอม 1 มีวิชาอะไรบ้าง''

\begin{verbatim}
SELECT apc.code, apc.name_th, apc.credits_total
FROM academic_plan_courses apc
JOIN programs p ON p.program_id = apc.program_id
WHERE p.program_code = 'IT'
  AND apc.plan_type = 'no_coop'
  AND apc.year = 1
  AND apc.semester = 1;
\end{verbatim}

\noindent\textbf{ตัวอย่างที่ 4} คำถาม: ``เกณฑ์การพ้นสภาพ (รีไทร์) เป็นอย่างไร''

\begin{verbatim}
SELECT category, text
FROM rules
WHERE category = 'dismissal_criteria'
  AND text LIKE '%เฉลี่ยสะสม%';
\end{verbatim}
